% Section: truth-in-model

\documentclass[../../../include/open-logic-section]{subfiles}

\begin{document}

\olfileid[pt]{nml}{syn}{tru}

\olsection{Verdade em um modelo}

Às vezes, interessa-nos saber quais !!{formula}s são verdadeiras em todos
os mundos de um dado modelo. Introduzamos uma notação para isso.

\begin{defn}
  !!^a{formula}~$!A$ é \emph{verdadeira em um modelo} $M = \tuple{W, R,
    V}$, o que se escreve $\mSat{M}{!A}$, se, e somente se,
  $\mSat{M}{!A}[w]$ para todo $w \in W$.
\end{defn}

\begin{prop}\ollabel{prop:truthfacts}
  \begin{enumerate}
  \item Se $\mSat{M}{!A}$, então $\mSat/{M}{\lnot !A}$, mas o inverso
    \emph{não} vale.
  \item Se $\mSat{M}{!A \lif !B}$, então $\mSat{M}{!A}$ somente se
    $\mSat{M}{!B}$, mas o inverso \emph{não} vale.
  \end{enumerate}
\end{prop}

\begin{proof}
  \begin{enumerate}
  \item Se $\mSat{M}{!A}$, então $!A$ é verdadeira em todos os mundos
    de $W$ e, como $W \neq \emptyset$, não pode ocorrer
    $\mSat{M}{\lnot!A}$; do contrário, $!A$ teria de ser verdadeira e
    falsa em algum mundo.

    Por outro lado, se $\mSat/{M}{\lnot !A}$, então $!A$ é verdadeira
    em algum mundo $w \in W$. Não se segue que $\mSat{M}{!A}[w]$ para
    \emph{todo}~$w \in W$. Por exemplo, no modelo de
    \olref[rel]{fig:simple}, $\mSat/{M}{\lnot p}$ e também $\mSat/{M}{p}$.
  \item Suponha $\mSat{M}{!A \lif !B}$ e $\mSat{M}{!A}$; para mostrar
    $\mSat{M}{!B}$, seja $w \in W$ um mundo arbitrário. Então
    $\mSat{M}{!A \lif !B}[w]$ e $\mSat{M}{!A}[w]$; logo,
    $\mSat{M}{!B}[w]$ e, como $w$ era arbitrário,
    $\mSat{M}{!B}$.

    Para mostrar que o inverso falha, precisamos encontrar um modelo
    $\mModel{M}$ tal que $\mSat{M}{!A}$ somente se $\mSat{M}{!B}$, mas
    $\mSat/{M}{!A \lif !B}$. Considere novamente o modelo de
    \olref[rel]{fig:simple}: $\mSat/{M}{p}$ e, portanto, vacuamente,
    $\mSat{M}{p}$ somente se $\mSat{M}{q}$. Contudo, $\mSat/{M}{p \lif
      q}$, pois $p$ é verdadeira, mas $q$ é falsa em~$w_1$.
  \end{enumerate}
\end{proof}

\begin{prob}
  Considere o modelo $\mModel{M}$ a seguir para a linguagem cujas únicas
  !!{propositional variable}s são $p_1$, $p_2$, $p_3$:
  \begin{center}
    \begin{tikzpicture}[modal]
      \node[world] (w1) [label={[align=right]left:\mTrue{p_1}\\\mFalse{p_2}\\\mFalse{p_3}}]
            {$w_1$} ; 
      \node[world] (w2) [label={[align=right]right:\mTrue{p_1}\\\mTrue{p_2}\\\mFalse{p_3}},
        below right=of w1]  {$w_2$}; 
      \node[world] (w3) [label={[align=right]right:\mTrue{p_1}\\\mTrue{p_2}\\\mTrue{p_3}},
        above right=of w2] {$w_3$};
      \draw[reflexive above] (w3) to (w3); 
      \draw[->] (w1) to (w2);
      \draw[->] (w2) to (w3);
      \draw[->] (w1) to (w3); 
    \end{tikzpicture}
  \end{center}
  As !!{formula}s e os esquemas a seguir são verdadeiros no modelo
  $\mModel{M}$, isto é, verdadeiros em todos os mundos de $\mModel{M}$?
  Explique.
  \begin{enumerate}
  \item $p\lif \Diamond p$ (para $p$ atômica);
  \item $!A\lif \Diamond !A$ (para $!A$ arbitrária);
  \item $\Box p \lif p$ (para $p$ atômica);
  \item $\lnot p \lif \Diamond \Box p$ (para $p$ atômica);
  \item $\Diamond \Box !A$ (para $!A$ arbitrária);
  \item $\Box \Diamond p$ (para $p$ atômica). 
  \end{enumerate}
\end{prob}

\end{document}
